\documentclass[a4paper,11pt]{article}
\usepackage[a4paper,bindingoffset=0in,left=0.9in,right=0.9in,top=1in,bottom=1in,footskip=.25in]{geometry}
\usepackage{generic}
\usepackage{comment}
\usepackage{graphicx}
\usepackage{epstopdf}
\usepackage{amsmath, amsthm}
\usepackage{setspace}
% Standard IEEE citation
\usepackage{cite}
 \usepackage{cuted}
%\usepackage[ruled,vlined]{algorithm2e}
 \usepackage{flushend}
 %\usepackage{multicol}
% Hyperlinks for references, URLs, sections
\usepackage[colorlinks=true, linkcolor=blue, citecolor=blue, urlcolor=blue]{hyperref}
\newtheorem{theorem}{Theorem}[section]
\newtheorem{proposition}{Proposition}[section]
\newtheorem{remark}{Remark}[section]
\newtheorem{lemma}{Lemma}[section]
% Math and symbols
\usepackage{amsmath,amssymb,amsfonts}
\usepackage{enumitem}
\setlist{noitemsep, topsep=0pt}

% Graphics
\usepackage{graphicx}
\setlength{\parindent}{0pt}
\usepackage[justification=centering]{caption}
% Algorithms
\usepackage{algorithm,algorithmic}
% \usepackage{enumerate}
\usepackage{rotating}
% Cross-referencing
\usepackage[capitalize]{cleveref}
\crefname{section}{Sec.}{Secs.}
\Crefname{section}{Section}{Sections}
\Crefname{table}{Table}{Tables}
\crefname{table}{Tab.}{Tabs.}

\def\BibTeX{{\rm B\kern-.05em{\sc i\kern-.025em b}\kern-.08em
    T\kern-.1667em\lower.7ex\hbox{E}\kern-.125emX}}


\title{Classical and Bayesian estimation of ....................... distribution}
\author{Student 1 \& Student 2 \\ Mehta Family School of Data Science and AI, IIT Guwahati, India }
\date{}
\begin{document}

\begin{singlespace}
 \maketitle
\end{singlespace}
% % \bibliographystyle{agsm}
% \setcounter{page}{1}
% \setcounter{equation}{0}
% \setcounter{figure}{0}
% \setcounter{table}{0}
% \newcommand{\oddhead}{Kumaraswamy}
% \renewcommand{\@oddhead}
% {\hspace*{-3pt}\raisebox{-3pt}[\headheight][0pt]
% 	{\vbox{\hbox to \textwidth
% 			{\hfill\oddhead}\vskip8pt}}}
% \vspace*{0.5in}
%\begin{onehalfspace}

\begin{abstract}
\noindent 


\noindent In this article, we consider the Bayesian estimation of ................. The performance of different estimators is evaluated through extensive numerical simulations in terms of mean-squared error and absolute bias, while ACI and HPD credible intervals are assessed in terms of average width and coverage probability. Finally, ............real data are analyzed for illustration purposes.\\
%\noindent {\bf Mathematics Subject Classification:}  62F10, 62F15, 62N02\\
\textbf{Keywords:} ..... distribution; ...... prior; maximum likelihood estimator; Bayesian estimation.\\
\end{abstract}

%\doublespacing
\maketitle
\section{Introduction}
\label{sec:introduction}
Reliability testing is one of the important methods to assess the failure pattern of different products. In such situations, items are put on test under some operational condition and then failure times are recorded. ..........................................................................
\section{Maximum Likelihood Estimation}
.........................
\subsection{MLE Under no order restriction}
.............................
\subsection{MLE Under order restriction}
.............................
\section{Bayesion Estimation}
.............................
\subsection{Bayesion Estimation under No Order Restriction}
.............................
\subsection{Bayesion Estimation under Order Restriction}
............................................................

The organization of the paper is as follows. The DE of a RA distribution under type-II censoring is determined in section \ref{sec2}. The maximum likelihood estimates of the entropy $H$ and the Fisher information matrix under the given censoring are provided in Section \ref{sec3}. In Section \ref{sec4}, objective Bayesian analysis under non-informative priors is examined. Section \ref{sec5} presents a simulation study and demonstrates the efficacy of entropy measure in RA distribution using two real data applications. Finally, the conclusion is given in Section \ref{sec6}. 




\section{Numerical Illustration}\label{sec5}
We now examine the performance of MLE and Bayes estimators of DE under the type-II censoring scheme, and  evaluate proposed estimates using average bias (AB) and average mean squared error (AMSE) value. The Bayesian estimation algorithm and simulation study are presented in following subsection.

\subsection{Simulation Study}
In this section, we conduct a simulation study to evaluate the performance of classical and Bayesian estimators ..................................


\subsection{Real Data Analysis}

In this section, we present an analysis of real data for illustrative purposes. 

%\newpage
\section{Conclusion} \label{sec6}
In this paper, we develop estimators of \\


%\newpage
%%%%%%%%%%%%%%%%%%%%%%%%%%%%
%\section{References}
\bibliographystyle{IEEEtran}
\bibliography{reference}

\end{document}